% !TEX program = xelatex
% ============================================================
%  Deep Generative Models — Chapter 2: Variational Autoencoder
% ============================================================

\documentclass[11pt,a4paper]{article}

% ============================================================
% §0  Fonts
% ============================================================
\usepackage{ctex}

% ============================================================
% §1  Page layout (A4, aligned with convex optimization notes)
% ============================================================
\IfFileExists{geometry.sty}{%
  \usepackage[margin=2.5cm]{geometry}
}{%
  \PackageWarning{geometry}{geometry package not found; using default A4 margins. Install `geometry' for 2.5cm margins.}%
}
\usepackage{microtype}
\linespread{1.08}

% ============================================================
% §2  Packages
% ============================================================
\usepackage{amsmath,amssymb,amsthm}
\renewcommand{\proofname}{Proof}
\usepackage{mathtools}
\usepackage{graphicx}
\graphicspath{{../graph/}}
\usepackage{float}
\usepackage{enumitem}
\usepackage{soul}
\usepackage{leftindex}
\usepackage{tikz}
\usetikzlibrary{arrows,calc,positioning}
\usepackage[dvipsnames]{xcolor}
\usepackage[most]{tcolorbox}
\usepackage{fontawesome5}
\usepackage{titlesec}

% ============================================================
% §3  Colors
% ============================================================
\definecolor{mydarkblue}{rgb}{0.0,0.08,0.45}
\definecolor{nicepink}  {RGB}{255,105,180}
\definecolor{cardbg}    {RGB}{245,245,247}
\definecolor{blue_new}  {RGB}{63,77,238}
\definecolor{page_white}{RGB}{227,227,229}
\definecolor{highlight} {HTML}{FCD66F}
\definecolor{warning}   {HTML}{E57373}
\definecolor{info}      {HTML}{4DB6AC}
\definecolor{navblue}     {HTML}{0077b6}
\definecolor{navcyan}     {HTML}{00b4d8}
\definecolor{navdark}     {HTML}{023e8a}
\definecolor{accentorange}{HTML}{fb8500}

\newcommand{\red}   [1]{{\color{OrangeRed}#1}}
\newcommand{\tlight}[1]{\textcolor{warning}{#1}}
\newcommand{\tinfo} [1]{\textcolor{info}{#1}}
\newcommand{\bglight}[1]{\hl{#1}}
\newcommand{\cbox}[2]{\colorbox{#1}{\textcolor{black}{#2}}}
\newcommand{\bgwarn}[1]{\cbox{warning}{#1}}
\newcommand{\bginfo}[1]{\cbox{info}{#1}}

% ============================================================
% §4  Hyperlinks
% ============================================================
\usepackage{hyperref}
\hypersetup{
    colorlinks = true,
    linkcolor  = mydarkblue,
    citecolor  = mydarkblue,
    urlcolor   = mydarkblue
}

% ============================================================
% §5  Math macros
% ============================================================
\newcommand{\wh}{\widehat}
\newcommand{\wt}{\widetilde}
\newcommand{\ov}{\overline}
\renewcommand{\tilde}{\wt}
\renewcommand{\hat}{\wh}
\renewcommand{\d}{\mathrm{d}}

\newcommand{\N}{\mathbb{N}}  \newcommand{\R}{\mathbb{R}}
\newcommand{\Q}{\mathbb{Q}}  \newcommand{\Z}{\mathbb{Z}}
\newcommand{\C}{\mathbb{C}}

\DeclareMathOperator*{\E}{\mathbb{E}}
\let\P\relax
\DeclareMathOperator*{\P}{\mathbb{P}}
\DeclareMathOperator*{\var}{Var}
\DeclareMathOperator*{\argmin}{arg\,min}
\DeclareMathOperator*{\argmax}{arg\,max}
\DeclareMathOperator{\diag}{diag}
\DeclareMathOperator{\tr}{tr}
\DeclareMathOperator{\rank}{rank}
\DeclareMathOperator{\dist}{dist}

\renewcommand{\L}{\mathcal{L}}

% ---- bold vectors ----
\newcommand{\ba}{\mathbf{a}} \newcommand{\bb}{\mathbf{b}}
\newcommand{\bc}{\mathbf{c}} \newcommand{\bd}{\mathbf{d}}
\newcommand{\be}{\mathbf{e}} \newcommand{\bg}{\mathbf{g}}
\newcommand{\bh}{\mathbf{h}} \newcommand{\bi}{\mathbf{i}}
\newcommand{\bj}{\mathbf{j}} \newcommand{\bk}{\mathbf{k}}
\newcommand{\bl}{\mathbf{l}} \newcommand{\bm}{\mathbf{m}}
\newcommand{\bn}{\mathbf{n}} \newcommand{\bo}{\mathbf{o}}
\newcommand{\bp}{\mathbf{p}} \newcommand{\bq}{\mathbf{q}}
\newcommand{\br}{\mathbf{r}} \newcommand{\bs}{\mathbf{s}}
\newcommand{\bt}{\mathbf{t}} \newcommand{\bu}{\mathbf{u}}
\newcommand{\bv}{\mathbf{v}} \newcommand{\bw}{\mathbf{w}}
\newcommand{\bx}{\mathbf{x}} \newcommand{\by}{\mathbf{y}}
\newcommand{\bz}{\mathbf{z}}
\newcommand{\bzero}{\mathbf{0}} \newcommand{\bI}{\mathbf{I}}
\newcommand{\bmu}{\boldsymbol{\mu}}
\newcommand{\bsigma}{\boldsymbol{\sigma}}
\newcommand{\beps}{\boldsymbol{\epsilon}}

% ---- bold matrices ----
\newcommand{\bA}{\mathbf{A}} \newcommand{\bB}{\mathbf{B}}
\newcommand{\bC}{\mathbf{C}} \newcommand{\bD}{\mathbf{D}}
\newcommand{\bE}{\mathbf{E}} \newcommand{\bF}{\mathbf{F}}
\newcommand{\bG}{\mathbf{G}} \newcommand{\bH}{\mathbf{H}}
\newcommand{\bJ}{\mathbf{J}} \newcommand{\bK}{\mathbf{K}}
\newcommand{\bL}{\mathbf{L}} \newcommand{\bM}{\mathbf{M}}
\newcommand{\bN}{\mathbf{N}} \newcommand{\bO}{\mathbf{O}}
\newcommand{\bP}{\mathbf{P}} \newcommand{\bQ}{\mathbf{Q}}
\newcommand{\bR}{\mathbf{R}} \newcommand{\bS}{\mathbf{S}}
\newcommand{\bT}{\mathbf{T}} \newcommand{\bU}{\mathbf{U}}
\newcommand{\bV}{\mathbf{V}} \newcommand{\bW}{\mathbf{W}}
\newcommand{\bX}{\mathbf{X}} \newcommand{\bY}{\mathbf{Y}}
\newcommand{\bZ}{\mathbf{Z}}

% ---- calligraphic ----
\newcommand{\cA}{\mathcal{A}} \newcommand{\cB}{\mathcal{B}}
\newcommand{\cC}{\mathcal{C}} \newcommand{\cD}{\mathcal{D}}
\newcommand{\cE}{\mathcal{E}} \newcommand{\cF}{\mathcal{F}}
\newcommand{\cG}{\mathcal{G}} \newcommand{\cH}{\mathcal{H}}
\newcommand{\cI}{\mathcal{I}} \newcommand{\cJ}{\mathcal{J}}
\newcommand{\cK}{\mathcal{K}} \newcommand{\cL}{\mathcal{L}}
\newcommand{\cM}{\mathcal{M}} \newcommand{\cN}{\mathcal{N}}
\newcommand{\cO}{\mathcal{O}} \newcommand{\cP}{\mathcal{P}}
\newcommand{\cQ}{\mathcal{Q}} \newcommand{\cR}{\mathcal{R}}
\newcommand{\cS}{\mathcal{S}} \newcommand{\cT}{\mathcal{T}}
\newcommand{\cU}{\mathcal{U}} \newcommand{\cV}{\mathcal{V}}
\newcommand{\cW}{\mathcal{W}} \newcommand{\cX}{\mathcal{X}}
\newcommand{\cY}{\mathcal{Y}} \newcommand{\cZ}{\mathcal{Z}}

% ---- VAE-specific ----
\newcommand{\KL}[2]{D_{\mathrm{KL}}\!\left(#1\,\|\,#2\right)}
\newcommand{\elbo}{\cL_{\mathrm{ELBO}}}
\newcommand{\elboHVAE}{\cL_{\mathrm{ELBO}}^{\mathrm{HVAE}}}
\newcommand{\pdata}{p_{\mathrm{data}}}
\newcommand{\pprior}{p_{\mathrm{prior}}}

% ============================================================
% §6  Section heading style
% ============================================================
\titleformat{\section}
  {\large\bfseries\color{mydarkblue}}{\thesection}{0.5em}{}[\titlerule]
\titlespacing*{\section}{0pt}{1.5ex}{1ex}

\titleformat{\subsection}
  {\normalsize\bfseries}{\thesubsection}{0.5em}{}
\titlespacing*{\subsection}{0pt}{1ex}{0.5ex}

% ============================================================
% §7  Compact title + horizontal TOC
% ============================================================
\makeatletter
\newcommand{\makecompacttitle}{%
    \begin{tcolorbox}[
        colback=mydarkblue!5, colframe=mydarkblue!10,
        arc=3pt, boxrule=0.5pt,
        left=5pt, right=5pt, top=5pt, bottom=5pt
    ]
        {\large\bfseries \@title}\hfill{\small \@author}\\[-2pt]
        {\tiny\color{gray} \@date \hfill Deep Generative Models Notes}
    \end{tcolorbox}
    \vspace{-1em}
}
\makeatother

\newtcolorbox{tabNavBoxSmall}[1][]{
    enhanced,
    colframe=white, colback=gray!5, coltext=black!80,
    fontupper=\sffamily\small\bfseries,
    arc=1.5mm, boxrule=0pt, bottomrule=2pt,
    left=2mm, right=2mm, top=1.5mm, bottom=1.5mm,
    #1
}

\newcommand{\horizontalTOC}{%
    \begin{center}
    \begin{tcbraster}[
        raster columns=3, raster width=\linewidth,
        raster column skip=1mm, raster row skip=1mm, raster equal height
    ]
        \begin{tabNavBoxSmall}\hyperref[sec:model-structure]{\ Model Structure}\end{tabNavBoxSmall}
        \begin{tabNavBoxSmall}\hyperref[sec:encoder]{\ Why the Encoder}\end{tabNavBoxSmall}
        \begin{tabNavBoxSmall}\hyperref[sec:elbo]{\ ELBO}\end{tabNavBoxSmall}
        \begin{tabNavBoxSmall}\hyperref[sec:info]{\ Info-Theoretic View}\end{tabNavBoxSmall}
        \begin{tabNavBoxSmall}\hyperref[sec:reparam]{\ Reparameterization}\end{tabNavBoxSmall}
        \begin{tabNavBoxSmall}\hyperref[sec:gaussian-vae]{\ Gaussian VAE}\end{tabNavBoxSmall}
        \begin{tabNavBoxSmall}\hyperref[sec:hvae]{\ Hierarchical VAE}\end{tabNavBoxSmall}
        \begin{tabNavBoxSmall}\hyperref[sec:hvae-elbo]{\ HVAE ELBO}\end{tabNavBoxSmall}
        \begin{tabNavBoxSmall}\hyperref[sec:pathologies]{\ Pathologies}\end{tabNavBoxSmall}
    \end{tcbraster}
    \vspace{-2.5pt}%
    \noindent\textcolor{gray!30}{\rule{\linewidth}{1pt}}
    \end{center}
    \vspace{1em}
}

\newcommand{\navbox}[2]{%
    \hyperlink{section.#1}{%
        \begin{tcolorbox}[
            width=\linewidth, enhanced, center upper,
            colback=white, colframe=mydarkblue!20,
            arc=2pt, boxrule=0.8pt,
            top=2pt, bottom=2pt, left=1pt, right=1pt,
            fontupper=\tiny\bfseries, colupper=mydarkblue
        ]#2\end{tcolorbox}%
    }%
}

% ============================================================
% §8  Theorem environments (tcolorbox)
% ============================================================
\tcbset{
  mytheo/.style={
    enhanced, breakable,
    colback=#1!5!white, colframe=#1!70!black,
    coltitle=white, fonttitle=\bfseries\small,
    boxrule=0.6pt, arc=3pt,
    left=5pt, right=5pt, top=10pt, bottom=5pt,
    attach boxed title to top left={yshift*=-\tcboxedtitleheight/2, xshift=2mm},
    boxed title style={
      colback=#1!80!black, colframe=#1!80!black,
      boxrule=0pt, arc=2pt
    }
  }
}

\newtcbtheorem[auto counter]            {theorem}    {Theorem}{mytheo=WildStrawberry}{thm}
\newtcbtheorem[use counter from=theorem]{definition} {Definition}{mytheo=NavyBlue  }{def}
\newtcbtheorem[use counter from=theorem]{example}    {Example}{mytheo=YellowOrange }{ex}
\newtcbtheorem[use counter from=theorem]{remark}     {Remark}{mytheo=Periwinkle    }{rem}
\newtcbtheorem[use counter from=theorem]{proposition}{Proposition}{mytheo=JungleGreen}{prop}
\newtcbtheorem[use counter from=theorem]{lemma}      {Lemma}{mytheo=ForestGreen  }{lem}
\newtcbtheorem[use counter from=theorem]{corollary}  {Corollary}{mytheo=BrickRed   }{cor}
\newtcbtheorem[use counter from=theorem]{fact}       {Fact}{mytheo=YellowGreen  }{fact}
\newtcbtheorem[use counter from=theorem]{assumption} {Assumption}{mytheo=Plum     }{as}

\tcolorboxenvironment{proof}{
  blanker, breakable, left=4pt, right=4pt, top=3pt, bottom=3pt,
  borderline west={2pt}{0pt}{gray!50}
}

% ============================================================
% §9  Body
% ============================================================
\title{Chapter 2: Variational Autoencoder}
\author{}
\date{\today}

\begin{document}

\makecompacttitle
\horizontalTOC

\begin{center}
  \includegraphics[width=\linewidth]{\detokenize{ChatGPT Image 2026年6月11日 18_12_18.png}}
\end{center}
\vspace{0.5em}

% ============================================================
\section{Model Structure}\label{sec:model-structure}
% ============================================================

\subsection{Encoder and Decoder}

A VAE consists of two parameterized networks: encoder $q_\phi(z\mid \bx)$ and decoder $p_\theta(\bx\mid z)$.

\begin{definition}{Variational Autoencoder}{vae}
A \textbf{variational autoencoder} is a latent-variable generative model specified by
\begin{itemize}[leftmargin=1.2em, itemsep=2pt]
    \item a fixed \textbf{prior} $p(z)$ over a low-dimensional latent variable $z\in\R^d$;
    \item a \textbf{decoder} (generative model) $p_\theta(\bx\mid z)$, a probabilistic map $\R^d\to\R^p$ parameterized by a neural network with weights $\theta$;
    \item an \textbf{encoder} (inference model) $q_\phi(z\mid \bx)$, a probabilistic map $\R^p\to\R^d$ parameterized by a neural network with weights $\phi$, used to approximate the intractable posterior $p_\theta(z\mid \bx)$.
\end{itemize}
\end{definition}

\medskip
Information flow:
\[
\bx \;\xrightarrow{\text{encoder}}\; q_\phi(z\mid \bx)
  \;\longrightarrow\; z
  \;\xrightarrow{\text{decoder}}\; p_\theta(\bx\mid z)
  \;\longrightarrow\; \bx'
\]

\begin{remark}{From Deterministic to Probabilistic Autoencoders}{ae-vs-vae}
A classical autoencoder learns a deterministic map $\bx\mapsto z\mapsto \bx'$ for reconstruction. A VAE instead makes the latent representation probabilistic and ties it to a prior $p(z)$, so new samples are generated by drawing $z\sim p(z)$ and decoding.
\end{remark}

\subsection{Variables}

\begin{itemize}[leftmargin=1.2em, itemsep=2pt]
    \item $\bx$: observed variable
    \item $z$: latent variable with prior $\pprior(z) = \cN(z;\,\mathbf{0},\,\mathbf{I})$
    \item Generation: sample $z\sim\pprior$, then $\bx\sim p_\theta(\bx\mid z)$
\end{itemize}

\subsection{Marginal Likelihood}

\[
    p_\theta(\bx) = \int p_\theta(\bx\mid z)\,p(z)\,dz
    = \E_{z\sim p(z)}\!\left[p_\theta(\bx\mid z)\right].
\]

This integral is intractable because $p_\theta(\bx\mid z)$ is a nonlinear neural network, making the integral impossible to evaluate analytically.

% ============================================================
\section{Why We Need the Encoder}\label{sec:encoder}
% ============================================================

By Bayes' rule:
\[
    p_\theta(z\mid \bx) = \frac{p_\theta(\bx\mid z)\,p(z)}{p_\theta(\bx)}
\]

Since $p_\theta(\bx)$ is intractable, so is the true posterior. We introduce a variational inference network:
\[
    q_\phi(z\mid \bx)\approx p_\theta(z\mid \bx)
\]

\subsection{An Importance-Sampling Viewpoint}

The encoder can be understood as a learned \textbf{proposal distribution} for importance sampling. Rewriting the marginal likelihood with an arbitrary $q_\phi(z\mid \bx)>0$,
\[
    p_\theta(\bx)
    = \int q_\phi(z\mid \bx)\,\frac{p_\theta(\bx\mid z)\,p(z)}{q_\phi(z\mid \bx)}\,dz
    = \E_{z\sim q_\phi(z\mid \bx)}\!\left[\frac{p_\theta(\bx\mid z)\,p(z)}{q_\phi(z\mid \bx)}\right].
\]
Prior sampling is inefficient because most sampled $z$ contribute little to explaining a fixed $\bx$. The ideal proposal is the true posterior $p_\theta(z\mid \bx)$: then the importance weight $p_\theta(\bx\mid z)\,p(z)/p_\theta(z\mid \bx) = p_\theta(\bx)$ is a constant, so the estimator has zero variance and is maximally efficient. Since the true posterior is intractable, we learn $q_\phi(z\mid \bx)$ as a tractable approximation.

% ============================================================
\section{Training Objective: ELBO}\label{sec:elbo}
% ============================================================

\subsection{Lower Bound}

\begin{definition}{Evidence Lower Bound (ELBO)}{elbo-def}
For any datapoint $\bx$ and any encoder $q_\phi(z\mid \bx)$ with the same support as the joint $p_\theta(\bx,z)$, the \textbf{evidence lower bound} is
\[
    \elbo(\theta,\phi;\,\bx)
    \;\triangleq\;
    \E_{z\sim q_\phi(z\mid \bx)}\!\left[\log\frac{p_\theta(\bx,z)}{q_\phi(z\mid \bx)}\right]
    =
    \E_{z\sim q_\phi}\!\left[\log p_\theta(\bx\mid z)\right]
    - \KL{q_\phi(z\mid \bx)}{p(z)} .
\]
The name reflects that the marginal log-likelihood $\log p_\theta(\bx)$ is often called the \emph{evidence}.
\end{definition}

\begin{theorem}{ELBO Lower Bound}{elbo}
\[
    \log p_\theta(\bx) \;\geq\; \elbo(\theta,\phi;\,\bx)
    \;\triangleq\;
    \underbrace{\E_{z\sim q_\phi}\!\left[\log p_\theta(\bx\mid z)\right]}_{\text{reconstruction}}
    \;-\;
    \underbrace{\KL{q_\phi(z\mid \bx)}{p(z)}}_{\text{latent KL}}
\]
\end{theorem}

\begin{proof}
Since $\log p_\theta(\bx)$ does not depend on $z$, we may take an expectation over $z\sim q_\phi(z\mid \bx)$, then multiply and divide by $q_\phi(z\mid \bx)$ inside the logarithm:
\begin{align*}
    \log p_\theta(\bx)
    &= \E_{z\sim q_\phi}\!\left[\log p_\theta(\bx)\right]
     = \E_{z\sim q_\phi}\!\left[\log \frac{p_\theta(\bx,z)}{p_\theta(z\mid \bx)}\right] \\
    &= \E_{z\sim q_\phi}\!\left[\log \frac{p_\theta(\bx,z)}{q_\phi(z\mid \bx)}\right]
     + \underbrace{\E_{z\sim q_\phi}\!\left[\log \frac{q_\phi(z\mid \bx)}{p_\theta(z\mid \bx)}\right]}_{=\,\KL{q_\phi(z\mid \bx)}{p_\theta(z\mid \bx)}\,\geq\,0}\\
    &\geq \E_{z\sim q_\phi}\!\left[\log\frac{p_\theta(\bx,z)}{q_\phi(z\mid \bx)}\right]
    = \E_z[\log p_\theta(\bx\mid z)] - \KL{q_\phi(z\mid \bx)}{p(z)},
\end{align*}
where the last line uses $p_\theta(\bx,z)=p_\theta(\bx\mid z)\,p(z)$. The same bound follows directly from Jensen's inequality applied to the concave $\log$:
\[
    \log p_\theta(\bx)
    = \log \E_{z\sim q_\phi}\!\left[\frac{p_\theta(\bx,z)}{q_\phi(z\mid \bx)}\right]
    \geq \E_{z\sim q_\phi}\!\left[\log\frac{p_\theta(\bx,z)}{q_\phi(z\mid \bx)}\right].
\]
\end{proof}

\begin{remark}{Tightness of the Bound}{elbo-gap}
The ELBO gap is
\[
    \log p_\theta(\bx) - \elbo(\theta,\phi;\,\bx)
    = \KL{q_\phi(z\mid \bx)}{p_\theta(z\mid \bx)} \;\geq\; 0 .
\]
The bound is tight iff $q_\phi(z\mid \bx)=p_\theta(z\mid \bx)$ almost everywhere.
\end{remark}

\subsection{Interpretation}

\begin{itemize}[leftmargin=1.2em, itemsep=2pt]
    \item \textbf{Reconstruction term}: encourages recovery of $\bx$ from $z$
    \item \textbf{KL term}: regularizes $q_\phi(z\mid \bx)$ toward the prior $\pprior$
\end{itemize}

% ============================================================
\section{Information-Theoretic View of ELBO}\label{sec:info}
% ============================================================

\subsection{Two Joint Distributions}

\begin{itemize}[leftmargin=1.2em, itemsep=2pt]
    \item \textbf{Generative joint}: $p_\theta(\bx,z)=p(z)\,p_\theta(\bx\mid z)$
    \item \textbf{Inference joint}: $q_\phi(\bx,z)=\pdata(\bx)\,q_\phi(z\mid \bx)$
\end{itemize}

\begin{proposition}{Modeling Error Bound}{info}
\[
    \KL{\pdata(\bx)}{p_\theta(\bx)} \;\leq\; \KL{q_\phi(\bx,z)}{p_\theta(\bx,z)}
\]
\end{proposition}

\begin{proof}
Factor both joints by the chain rule, $q_\phi(\bx,z)=\pdata(\bx)\,q_\phi(z\mid \bx)$ and $p_\theta(\bx,z)=p_\theta(\bx)\,p_\theta(z\mid \bx)$, and split the logarithm of the ratio:
\begin{align*}
    \KL{q_\phi(\bx,z)}{p_\theta(\bx,z)}
    &= \E_{q_\phi(\bx,z)}\!\left[\log\frac{\pdata(\bx)\,q_\phi(z\mid \bx)}{p_\theta(\bx)\,p_\theta(z\mid \bx)}\right]\\
    &= \E_{q_\phi(\bx,z)}\!\left[\log\frac{\pdata(\bx)}{p_\theta(\bx)}\right]
     + \E_{q_\phi(\bx,z)}\!\left[\log\frac{q_\phi(z\mid \bx)}{p_\theta(z\mid \bx)}\right].
\end{align*}
The first term depends only on $\bx$, whose marginal under $q_\phi(\bx,z)$ is $\pdata(\bx)$, so it equals $\KL{\pdata}{p_\theta}$. In the second term, taking the inner expectation over $z\sim q_\phi(z\mid \bx)$ first turns each fiber into a KL divergence, leaving an outer expectation over $\bx\sim\pdata$:
\begin{align*}
    \KL{q_\phi(\bx,z)}{p_\theta(\bx,z)}
    &= \KL{\pdata}{p_\theta}
      + \E_{\pdata}\!\left[\KL{q_\phi(z\mid \bx)}{p_\theta(z\mid \bx)}\right]\\
    &\;\geq\;\KL{\pdata}{p_\theta},
\end{align*}
since the second term is an expectation of nonnegative KL divergences.
\end{proof}

\subsection{Error Decomposition}

\begin{itemize}[leftmargin=1.2em, itemsep=2pt]
    \item $\KL{\pdata}{p_\theta}$: \textbf{modeling error}
    \item $\E_{\pdata}\!\left[\KL{q_\phi(z\mid \bx)}{p_\theta(z\mid \bx)}\right]$: \textbf{inference error}
\end{itemize}

\begin{remark}{ELBO and Inference Error}{}
\[
    \log p_\theta(\bx) - \elbo(\theta,\phi;\,\bx)
    = \KL{q_\phi(z\mid \bx)}{p_\theta(z\mid \bx)} \geq 0
\]
Maximizing ELBO therefore reduces inference error.
\end{remark}

% ============================================================
\section{Reparameterization and Gradient Estimation}\label{sec:reparam}
% ============================================================

\subsection{The Gradient Problem}

Training maximizes $\E_{\pdata(\bx)}[\elbo(\theta,\phi;\bx)]$ by stochastic gradient ascent. The gradient with respect to the decoder $\theta$ is straightforward, since the sampling distribution $q_\phi$ does not depend on $\theta$:
\[
    \nabla_\theta\,\E_{z\sim q_\phi}\!\left[\log p_\theta(\bx\mid z)\right]
    = \E_{z\sim q_\phi}\!\left[\nabla_\theta\log p_\theta(\bx\mid z)\right].
\]
The encoder gradient $\nabla_\phi$ is harder because the sampling distribution itself depends on $\phi$:
\[
    \nabla_\phi\,\E_{z\sim q_\phi(z\mid \bx)}\!\left[f(z)\right]
    \neq \E_{z\sim q_\phi(z\mid \bx)}\!\left[\nabla_\phi f(z)\right].
\]

\subsection{The Reparameterization Trick}

\begin{definition}{Reparameterization}{reparam-def}
A distribution $q_\phi(z\mid \bx)$ is \textbf{reparameterizable} if there exist a fixed base distribution $p(\beps)$ (independent of $\phi$) and a differentiable map $g_\phi(\beps,\bx)$ such that
\[
    z\sim q_\phi(z\mid \bx)
    \quad\Longleftrightarrow\quad
    z = g_\phi(\beps,\bx),\;\; \beps\sim p(\beps).
\]
For the diagonal Gaussian encoder $q_\phi(z\mid \bx)=\cN\!\left(z;\bmu_\phi(\bx),\diag(\bsigma_\phi^2(\bx))\right)$ this is
\[
    z = \bmu_\phi(\bx) + \bsigma_\phi(\bx)\odot\beps,
    \qquad \beps\sim\cN(\bzero,\bI),
\]
where $\odot$ denotes the elementwise product.
\end{definition}

\begin{proposition}{Pathwise Gradient Estimator}{reparam-grad}
If $q_\phi(z\mid \bx)$ is reparameterizable via $z=g_\phi(\beps,\bx)$, then for any differentiable $f$,
\[
    \nabla_\phi\,\E_{z\sim q_\phi(z\mid \bx)}\!\left[f(z)\right]
    = \E_{\beps\sim p(\beps)}\!\left[\nabla_\phi\, f\!\left(g_\phi(\beps,\bx)\right)\right].
\]
\end{proposition}

\begin{proof}
Because $p(\beps)$ does not depend on $\phi$, the expectation operator and the gradient commute under the change of variables $z=g_\phi(\beps,\bx)$:
\[
    \nabla_\phi\,\E_{z\sim q_\phi(z\mid \bx)}\!\left[f(z)\right]
    = \nabla_\phi\,\E_{\beps\sim p(\beps)}\!\left[f\!\left(g_\phi(\beps,\bx)\right)\right]
    = \E_{\beps\sim p(\beps)}\!\left[\nabla_\phi\, f\!\left(g_\phi(\beps,\bx)\right)\right],
\]
where the last step exchanges $\nabla_\phi$ with the (now $\phi$-free) expectation. The chain rule $\nabla_\phi f = (\partial_z f)\,(\partial_\phi g_\phi)$ then propagates the gradient through the sampled $z$.
\end{proof}

\subsection{Monte Carlo ELBO and Its Gradient}

Combining reparameterization with a single Monte Carlo sample per datapoint gives the unbiased estimator actually used in practice. With $z=g_\phi(\beps,\bx)$ and $\beps\sim p(\beps)$,
\[
    \wh{\elbo}(\theta,\phi;\bx)
    = \log p_\theta\!\left(\bx\mid g_\phi(\beps,\bx)\right)
      - \KL{q_\phi(z\mid \bx)}{p(z)},
\]
and its gradient (when the KL term is available in closed form, see Section~\ref{sec:gaussian-vae}) is
\[
    \nabla_{\theta,\phi}\,\wh{\elbo}(\theta,\phi;\bx)
    = \nabla_{\theta,\phi}\!\left[
        \log p_\theta\!\left(\bx\mid g_\phi(\beps,\bx)\right)
        - \KL{q_\phi(z\mid \bx)}{p(z)}
    \right],
\]
which is an unbiased estimator of $\nabla_{\theta,\phi}\,\elbo(\theta,\phi;\bx)$ and is computed by ordinary backpropagation.

\begin{remark}{Why Not the Score-Function Estimator}{score-vs-reparam}
The general-purpose score-function (REINFORCE) estimator
\[
    \nabla_\phi\,\E_{z\sim q_\phi}\!\left[f(z)\right]
    = \E_{z\sim q_\phi}\!\left[f(z)\,\nabla_\phi\log q_\phi(z\mid \bx)\right]
\]
is unbiased and does not require reparameterization, but often has much higher variance. For continuous latents, pathwise gradients are typically more stable in practice.
\end{remark}

% ============================================================
\section{Gaussian VAE}\label{sec:gaussian-vae}
% ============================================================

\subsection{Parameterization}

\textbf{Encoder} (diagonal Gaussian):
\[
    q_\phi(z\mid \bx) \triangleq \cN\!\left(z;\;\bmu_\phi(\bx),\;\diag\!\left(\bsigma^2_\phi(\bx)\right)\right)
\]
where $\bmu_\phi:\R^p\to\R^d$, $\bsigma_\phi:\R^p\to\R^d_+$.

\medskip
\textbf{Decoder} (isotropic Gaussian):
\[
    p_\theta(\bx\mid z) \triangleq \cN\!\left(\bx;\;\mathbf{M}_\theta(z),\;\sigma^2\mathbf{I}\right)
\]
where $\mathbf{M}_\theta:\R^d\to\R^p$ and $\sigma>0$ is a fixed small constant.

\subsection{Equivalent Optimization}

Expanding the reconstruction term:
\[
    \E_z[\log p_\theta(\bx\mid z)]
    = -\frac{1}{2\sigma^2}\E_z\!\left[\|\bx-\mathbf{M}_\theta(z)\|^2\right]+\mathrm{const}
\]

So maximizing ELBO is equivalent to:
\[
    \min_{\theta,\phi}\;
    \E_{\pdata}\!\left[
        \E_{q_\phi}\!\left[\frac{1}{2\sigma^2}\|\bx-\mathbf{M}_\theta(z)\|^2\right]
        + \KL{q_\phi(z\mid \bx)}{\pprior(z)}
    \right]
\]
Here $\sigma$ controls the reconstruction--regularization trade-off. The reconstruction term is estimated via reparameterized sampling (Section~\ref{sec:reparam}), while the KL term is closed form.

\begin{proposition}{Closed-Form Gaussian KL}{gauss-kl}
When both $q_\phi$ and $\pprior$ are Gaussian with diagonal covariance,
\[
    \KL{\cN(\bmu,\diag(\bsigma^2))}{\cN(\mathbf{0},\mathbf{I})}
    = \frac{1}{2}\sum_{j=1}^{d}\!\left(\mu_j^2+\sigma_j^2-\log\sigma_j^2-1\right).
\]
\end{proposition}

\begin{proof}
Because both distributions factorize over the $d$ coordinates, the KL divergence is the sum of the per-coordinate divergences; it suffices to prove the one-dimensional case for $q=\cN(\mu,\sigma^2)$ and $p=\cN(0,1)$. Writing out the densities,
\[
    \log\frac{q(z)}{p(z)}
    = \log\frac{1}{\sigma} - \frac{(z-\mu)^2}{2\sigma^2} + \frac{z^2}{2}
    = -\frac{1}{2}\log\sigma^2 - \frac{(z-\mu)^2}{2\sigma^2} + \frac{z^2}{2}.
\]
Taking the expectation over $z\sim q$ and using $\E_q[(z-\mu)^2]=\sigma^2$ together with $\E_q[z^2]=\mu^2+\sigma^2$,
\[
    \KL{q}{p}
    = \E_q\!\left[\log\frac{q(z)}{p(z)}\right]
    = -\frac{1}{2}\log\sigma^2 - \frac{1}{2} + \frac{\mu^2+\sigma^2}{2}
    = \frac{1}{2}\!\left(\mu^2+\sigma^2-\log\sigma^2-1\right).
\]
Summing over the $d$ independent coordinates yields the stated formula.
\end{proof}

\begin{remark}{A Fully Differentiable Objective}{gauss-objective}
With reparameterized reconstruction and closed-form KL (Proposition~\ref{prop:gauss-kl}), the per-sample objective is fully differentiable in $(\theta,\phi)$ and optimized by standard backpropagation.
\end{remark}

% ============================================================
\section{Hierarchical VAE (HVAE)}\label{sec:hvae}
% ============================================================

\subsection{Motivation}

Hierarchical VAEs introduce a hierarchy of latent variables, allowing the model to capture features at multiple abstraction levels.

\subsection{Hierarchical Encoder-Decoder Chain}

\[
\bx \xrightarrow[q_\phi(z_1\mid \bx)]{p_\theta(\bx\mid z_1)} z_1
\xrightarrow[q_\phi(z_2\mid z_1)]{p_\theta(z_1\mid z_2)}
\cdots
\xrightarrow[q_\phi(z_L\mid z_{L-1})]{p_\theta(z_{L-1}\mid z_L)} z_L
\]

\subsection{Modeling}

The generative joint distribution is:
\[
    p_\theta(\bx,z_{1:L})
    = p_\theta(\bx\mid z_1)\prod_{i=2}^{L} p_\theta(z_{i-1}\mid z_i)\,p(z_L)
\]

The marginal data distribution becomes:
\[
    p_{\mathrm{HVAE}}(\bx)=\int p_\theta(\bx,z_{1:L})\,dz_{1:L}
\]

\subsection{Generation and Encoding}

\textbf{Generation} samples top latent $z_L\sim p(z_L)$, then decodes downward layer by layer to $z_1$, and finally generates $\bx$ from $p_\theta(\bx\mid z_1)$.

\medskip
\textbf{Encoding} follows a Markov factorization:
\[
    q_\phi(z_{1:L}\mid \bx)=q_\phi(z_1\mid \bx)\prod_{i=2}^{L} q_\phi(z_i\mid z_{i-1})
\]
This bottom-up Markov chain mirrors the top-down decoder and keeps inference/sampling tractable through simple conditional factors.

% ============================================================
\section{HVAE ELBO}\label{sec:hvae-elbo}
% ============================================================

\subsection{Derivation}

\begin{align*}
    \log p_{\mathrm{HVAE}}(\bx)
    &= \log \int p_\theta(\bx,z_{1:L})\,dz_{1:L} \\
    &= \log \int q_\phi(z_{1:L}\mid \bx)\,
       \frac{p_\theta(\bx,z_{1:L})}{q_\phi(z_{1:L}\mid \bx)}\,dz_{1:L} \\
    &= \log \E_{q_\phi(z_{1:L}\mid \bx)}
       \left[\frac{p_\theta(\bx,z_{1:L})}{q_\phi(z_{1:L}\mid \bx)}\right] \\
    &\ge \E_{q_\phi(z_{1:L}\mid \bx)}
       \left[\log\frac{p_\theta(\bx,z_{1:L})}{q_\phi(z_{1:L}\mid \bx)}\right]
       \;\triangleq\; \elboHVAE(\theta,\phi;\bx)
\end{align*}

\subsection{Three-Term Decomposition}

Substituting the generative and inference factorizations into $\log\frac{p_\theta(\bx,z_{1:L})}{q_\phi(z_{1:L}\mid \bx)}$ and using $z_0\equiv\bx$ yields a reconstruction term, adjacent-layer KL matching terms, and a top-level KL to the prior:

\begin{align*}
    \elboHVAE(\theta,\phi;\bx)
    &= \underbrace{
        \E_{q_\phi(z_1\mid \bx)}\!\left[\log p_\theta(\bx\mid z_1)\right]
       }_{\text{reconstruction term}} \\
    &\quad - \underbrace{
        \sum_{i=1}^{L-1}
        \E_{q_\phi(z_{1:L}\mid \bx)}
        \!\left[
            \KL{q_\phi(z_i\mid z_{i-1})}{p_\theta(z_i\mid z_{i+1})}
        \right]
       }_{\text{adjacent-layer matching terms}} \\
    &\quad - \underbrace{
        \E_{q_\phi(z_{L-1}\mid \bx)}
        \!\left[\KL{q_\phi(z_L\mid z_{L-1})}{p(z_L)}\right]
       }_{\text{top-level KL regularizer}}
\end{align*}
where $z_0\equiv \bx$ for notational convenience.

% ============================================================
\section{Pathologies of Deep Networks in Flat VAE}\label{sec:pathologies}
% ============================================================

\subsection{Reason 1: Deeper Networks Do Not Expand the Variational Family}

In the standard variational family
\[
    q_\phi(z\mid \bx)=\cN\!\left(z;\mu_\phi(\bx),\diag(\sigma_\phi^2(\bx))\right),
\]
deeper encoder networks may improve parameter fitting quality, but they do not change the family itself. If the true posterior $p_\theta(z\mid \bx)$ is multi-modal, a diagonal Gaussian approximation remains fundamentally limited.

\subsection{Reason 2: Posterior Collapse}

The expected ELBO over data can be decomposed as:
\[
    \E_{\pdata(\bx)}[\elbo(\bx)]
    =
    \E_{\pdata(\bx)\,q_\phi(z\mid \bx)}[\log p_\theta(\bx\mid z)]
    - I_\phi(\bx;z)
    - \KL{\bar q_\phi(z)}{p(z)},
\]
where the aggregated posterior is
\[
    \bar q_\phi(z) \triangleq \E_{\pdata(\bx)}[q_\phi(z\mid \bx)],
\]
and the mutual information is
\[
    I_\phi(\bx;z)
    \triangleq
    \E_{\pdata(\bx)}
    \!\left[\KL{q_\phi(z\mid \bx)}{\bar q_\phi(z)}\right].
\]

This decomposition follows by averaging the per-datapoint latent KL and inserting the aggregated posterior $\bar q_\phi(z)$ inside the logarithm:
\begin{align*}
    \E_{\pdata(\bx)}\!\left[\KL{q_\phi(z\mid \bx)}{p(z)}\right]
    &= \E_{\pdata(\bx)\,q_\phi(z\mid \bx)}\!\left[\log\frac{q_\phi(z\mid \bx)}{\bar q_\phi(z)} + \log\frac{\bar q_\phi(z)}{p(z)}\right]\\
    &= \underbrace{\E_{\pdata(\bx)}\!\left[\KL{q_\phi(z\mid \bx)}{\bar q_\phi(z)}\right]}_{=\,I_\phi(\bx;z)}
     + \underbrace{\E_{\bar q_\phi(z)}\!\left[\log\frac{\bar q_\phi(z)}{p(z)}\right]}_{=\,\KL{\bar q_\phi(z)}{p(z)}},
\end{align*}
where the second term uses that the marginal of $z$ under $\pdata(\bx)\,q_\phi(z\mid \bx)$ is exactly $\bar q_\phi(z)$. Substituting back into $\E_{\pdata}[\elbo]=\E_{\pdata\,q_\phi}[\log p_\theta(\bx\mid z)]-\E_{\pdata}[\KL{q_\phi(z\mid \bx)}{p(z)}]$ gives the three-term form above.

If the decoder is too expressive, the model class may contain solutions with
$p_\theta(\bx\mid z)\approx r(\bx)\approx \pdata(\bx)$ that almost ignore $z$. Then $q_\phi(z\mid \bx)$ is pushed toward $p(z)$, implying $I_\phi(\bx;z)\to 0$ and KL penalties close to zero.

\begin{remark}{Consequences of Posterior Collapse}{}
\begin{itemize}[leftmargin=1.2em, itemsep=2pt]
    \item $z$ becomes nearly independent of $\bx$, carrying little useful information.
    \item Conditioning on $z$ or moving in latent space has little effect on generated samples.
\end{itemize}
\end{remark}

\subsection{What Hierarchy Changes?}

One structural response to posterior collapse is to adopt hierarchical latents (HVAE; Section~\ref{sec:hvae}). Instead of routing all information through a single code, the model distributes representation across multiple levels, with adjacent-layer matching terms that constrain inference and generation at every interface.

The HVAE ELBO (Section~\ref{sec:hvae-elbo}) can be written in log-difference form as
\begin{align*}
    \elboHVAE(\theta,\phi;\bx)
    &= \E_{q_\phi(z_1\mid \bx)}\!\left[\log p_\theta(\bx\mid z_1)\right] \\
    &\quad - \E_{q_\phi(z_{1:2}\mid \bx)}
       \!\left[\log q_\phi(z_1\mid \bx) - \log p_\theta(z_1\mid z_2)\right] \\
    &\quad - \sum_{i=2}^{L-1}
       \E_{q_\phi(z_{1:i+1}\mid \bx)}
       \!\left[\log q_\phi(z_i\mid z_{i-1}) - \log p_\theta(z_i\mid z_{i+1})\right] \\
    &\quad - \E_{q_\phi(z_{L-1}\mid \bx)}
       \!\left[\KL{q_\phi(z_L\mid z_{L-1})}{p(z_L)}\right].
\end{align*}
Under $z_0\equiv\bx$, each log-difference term is exactly the corresponding adjacent-layer KL matching term in Section~\ref{sec:hvae-elbo}; the top-level term is already a KL.

\begin{remark}{Layer-Wise Abstraction in HVAE}{}
Each latent level interacts only with its neighbors: the encoder passes information upward through $q_\phi(z_i\mid z_{i-1})$, and the decoder passes information downward through $p_\theta(z_{i-1}\mid z_i)$. These properties stem from the hierarchical latent graph, not from simply deepening networks in a flat VAE. Because matching penalties act at every layer, a powerful decoder cannot ignore the entire latent hierarchy without incurring multiple interface costs; HVAE is therefore typically less collapse-prone than a flat VAE with a deeper encoder alone.
\end{remark}

\end{document}
